\section{Numerical output stability without an all-input margin}
\label{sec:numerical-stability}

The preceding sign-robustness conditions can be replaced by a condition on
\emph{numerical changes in the output}. This permits the clean output to be
arbitrarily close to zero on unobserved inputs. Only the observed examples
are required to have the original positive training margin.

Let $J=L-1$, $E_v(w)=N\|v\|_2^2$, and set
\begin{equation}\label{eq:numeric-radius}
 r_{\rm num}(w)=\frac{1}{1000\sqrt{1+J E_v(w)}}.
\end{equation}
Use the same additive postactivation and scalar output disturbances as
before. Say that $w$ is \emph{numerically stable} if
\begin{equation}\label{eq:numeric-support}
 |z_w(x;\eta)-z_w(x)|\le\frac1{16}
 \quad\text{for all }x\in\{-1,1\}^d
 \text{ and all }\|\eta\|_*\le r_{\rm num}(w).
\end{equation}
The prior is the raw Gaussian conditioned on this event. There is
\emph{no all-input clean-margin condition}. For an observed labeled sample,
condition further on
\begin{equation}\label{eq:numeric-likelihood}
 y_i z_w(x_i)>\frac12\qquad(1\le i\le n).
\end{equation}
Thus the original output-margin convention is imposed only at training
inputs. With no global margin restriction, this likelihood is stronger
than sign-only consistency; that distinction is part of the theorem.

The condition has a simple sufficient example. If every hidden-to-hidden
matrix has operator norm at most one, then
\[
 |z_w(x;\eta)-z_w(x)|
 \le\sqrt{1+J E_v(w)}\,\|\eta\|_*.
\]
Consequently every such network satisfies~\eqref{eq:numeric-support},
regardless of how small its clean outputs are. The first layer is
unrestricted because there is no input perturbation, and the final readout
may have arbitrary gain. Numerical stability can also hold through
saturation and cancellation without nonexpansiveness.

\begin{lemma}[Numerical posterior compression]
\label{lem:numeric-compression}
Let $0<r\le1/6$, and suppose
$D_{\rm KL}(Q\Vert\mu_{N,t})\le NK/t^2$. Suppose every network in the
support of $Q$ obeys
\[
 |z_w(x;\eta)-z_w(x)|\le1/16\qquad(\|\eta\|_*\le r)
\]
at every input $x$ in a specified set $G$. There is a finite class
$\mathcal H$ of \emph{real-valued} depth-$L$ network outputs with
\[
 \log|\mathcal H|\le\operatorname{poly}(d,L,K,r^{-1},\alpha^{-1}),
\]
independent of $N,t,Q$, and a law $\Gamma_Q$ on $\mathcal H$, such that
at each $x\in G$ a coupling with marginals $Q,\Gamma_Q$ satisfies
\[
 \mathbb P\{|h(x)-z_w(x)|>1/12\}\le\alpha.
\]
The candidate marginal is the same at every input.
\end{lemma}
\begin{proof}
Run the earlier entropy compressor before using its final sign-robustness
step. Its entropy bounds and construction do not require robustness.
Except with probability less than $\alpha/4$, its unrounded candidate
output at a fixed input is exactly the original network's output under
a simultaneous postactivation disturbance of norm less than $r/2$.
The numerical stability assumption bounds the change by $1/16$.
The same fixed parameter-grid rounding changes outputs by at most $r/8$,
uniformly on the cube. Since $r\le1/6$,
$1/16+r/8\le1/12$. In the small-width branch the unrounded candidate is
the original network except on its bounded-weight failure event, and the
same rounding bound suffices. Keep the rounded real outputs, rather than
only their signs. The finite grid and the candidate marginal are otherwise
unchanged. No clean output margin has been used.
\end{proof}

\begin{theorem}[Polynomial-depth learning without a global clean margin]
\label{thm:numeric-prior}
For the prior~\eqref{eq:numeric-support} and training
likelihood~\eqref{eq:numeric-likelihood}, a target computed by a width-$C$
layered Boolean circuit has a sample bound
\[
 n\ge\operatorname{poly}(d,L,C,\varepsilon^{-1},\log(1/\delta))
\]
such that, with probability at least $1-\delta$,
\begin{equation}\label{eq:numeric-risk}
 \mathbb E_{w\sim Q_S,\,X\sim D}
       \mathbf1\{f(X)z_w(X)\le1/4\}\le\varepsilon.
\end{equation}
One sample event gives this guarantee simultaneously for every $N\ge C$
and $0<t\le1$, hence for every $\gamma\ge1$. The polynomial's degree is
independent of depth, width, and Gaussian scale. In particular, it bounds
Gibbs sign error, and gives the usual factors three and two for posterior
confidence thresholds $2/3$ and $1/2$, respectively.
\end{theorem}

\begin{proof}[Statistical proof, given the target cylinder below]
The target cylinder established below has raw Gaussian mass at least
$\exp(-ND_*/t^2)$ with $D_*=\operatorname{poly}(d,L,C)$, is numerically
stable, and has correct signed clean output greater than one on the full
cube. It therefore lies in the posterior conditioning event for every
sample. As before, posterior KL is at most $ND_*/t^2$.
The Gaussian readout tail and this partition lower bound allow discarding
posterior mass at most $\beta$ while retaining $E_v\le B$, where
$B=\operatorname{poly}(D_*,\log(1/\beta))$.
The retained law $Q'$ has KL at most $N(D_*+1)/t^2$ and numerical stability
at the common radius
$r_0=[1000\sqrt{1+JB}]^{-1}<1/6$.

Apply Lemma~\ref{lem:numeric-compression}, and define the bounded ramp
\[
 \psi(u)=
 \begin{cases}
 1,&u\le1/3,\\
 5-12u,&1/3<u<5/12,\\
 0,&u\ge5/12.
 \end{cases}
\]
On a successful comparison at a training point, the candidate has
$y_i h(x_i)>1/2-1/12=5/12$, so its ramp loss is zero. Thus the empirical
ramp risk averaged under $\Gamma_{Q'}$ is at most $\alpha$.
On a successful comparison at a test point, if $f(x)z_w(x)\le1/4$, then
$f(x)h(x)\le1/4+1/12=1/3$, so the candidate's ramp loss is one.
Therefore the original test margin-failure probability is at most the
candidate's population ramp risk plus $\alpha$.

The finite-class relative Chernoff inequality applies to these $[0,1]$
losses, and gives simultaneously for all $h\in\mathcal H$
\[
 R_\psi(h)\le2\widehat R_\psi(h)
             +\frac{8[\log|\mathcal H|+\log(1/\delta)]}{n}.
\]
For completeness, the bounded-loss version follows from the same
exponential argument as for indicators, using
$e^{-s u}\le1-(1-e^{-s})u$ for $0\le u\le1$.
Averaging under the candidate law and adding the discarded posterior mass
gives
\[
 \mathbb E_{Q_S,D}\mathbf1\{f(X)z_w(X)\le1/4\}
 \le\beta+3\alpha+
       \frac{8[\log|\mathcal H|+\log(1/\delta)]}{n}.
\]
Choose $\beta=\varepsilon/4$, $\alpha=\varepsilon/12$, and $n$ large enough
that the last term is at most $\varepsilon/2$.
The class is independent of $N,t$, proving simultaneity.
Finally, if the posterior gives probability at most $2/3$ to a correct
signed output greater than $1/4$, its conditional margin-failure probability
is at least $1/3$; Markov gives the factor three. The sign event is larger,
and the corresponding factor at $1/2$ is two.
\end{proof}

\subsection{A numerically stable target cylinder}

The robust digital-cylinder proof from the preceding report can be made
numerically stable by tightening its hidden approximation constants. Use
balanced clone groups, write $P$ for the group-constant projection, and
write $b^s(x)\in\{-1,1\}^N$ for the ideal repeated Boolean wires. Replace
that construction's parameters by
\begin{align*}
 e_*&=\frac1{128\sqrt{2C}},&
 \tau_0&=e_*/4,\\
 b_L&=\sqrt{16\log9+8\log(4L)},&
 A&=\max\{32b_L,\,2\log(8/e_*)\},\\
 \zeta&=\min\{1/4,\,A/[4(d+1)],\,Ae_*/(64\sqrt C)\}.
\end{align*}
Here $\zeta$ is the cylinder tolerance, unrelated to an input-distribution
exception mass. Constrain the standardized first-layer coordinates, active
internal group coordinates, and readout coordinates within $\zeta t$ of
the cloned saturated circuit, exactly as in that construction. On its
free-bulk event,
\[
 W^s=W_*^s+R^s+E^s,\quad
 R^s=R^sP,\quad \|R^s\|_{\rm op}\le\zeta t\sqrt C,\quad
 E^sP=0,\quad \|E^s\|_{\rm op}\le b_Lt.
\]
The free-bulk event has conditional probability at least $1/2$.

The old induction, with these tightened constants, proves
\begin{equation}\label{eq:numeric-tube}
 \frac{\|h^s_\eta(x)-b^s(x)\|_2}{\sqrt N}\le e_*
\end{equation}
for every input, at every hidden layer, whenever each hidden injection has
RMS at most $\tau_0$ and output noise has magnitude at most $\tau_0$.
To check the constants directly, an error of RMS $e_*$ makes every clone
group's mean error at most $\sqrt{2C}e_*=1/128$. Thus every ideal gate's
signed preactivation exceeds $A/2$. The global secant estimate
\[
 |\tanh(a+u)-b|\le2e^{-A/2}+(8/A)|u|
 \quad\text{when }b\in\{-1,1\},\ ba\ge A/2
\]
gives the recurrence
\[
 e_s\le2e^{-A/2}+(8/A)\zeta\sqrt C(1+\tau_0)
               +(8b_L/A)e_{s-1}+\tau_0\le e_*.
\]
Each of the four terms is at most $e_*/4$.
The first-layer perturbation changes each preactivation by at most
$\zeta(d+1)\le A/4$, so its signed preactivation is at least $3A/4$;
its activation error plus injection is at most $e_*/2$.
This initializes the recurrence. In particular the clean computation
also satisfies~\eqref{eq:numeric-tube}.

The cloned readout obeys $\sqrt N\|v_*\|\le\sqrt{8C}$, while its
parameter perturbation satisfies $\sqrt N\|v-v_*\|\le\zeta$.
Comparing the noisy and clean computations for the \emph{same} weights
therefore yields
\begin{align*}
 |z_w(x;\eta)-z_w(x)|
 &\le2e_*(\sqrt{8C}+\zeta)+\tau_0\\
 &\le\frac1{32}+\frac34 e_*<\frac1{16}.
\end{align*}
Also $\sqrt N\|v_*\|\ge\sqrt{2C}$, so
\[
 E_v(w)\ge(\sqrt{2C}-1/4)^2\ge C.
\]
Consequently
\[
 r_{\rm num}(w)\le\frac{1}{1000\sqrt C}
       <\frac1{512\sqrt{2C}}=\tau_0.
\]
Thus every draw in the cylinder and free-bulk event satisfies the numerical
support rule. Its correct signed clean output is at least
$2-\sqrt{8C}e_*-\zeta>1$, which also establishes the training likelihood.

There are still only $N[d+2+(L-2)C]$ constrained Gaussian coordinates,
and the standardized center energy is at most
$N[9dA^2+18(L-2)CA^2+8C]$. The preceding Gaussian box calculation therefore
gives mass at least $\exp(-ND_*/t^2)$ for
$D_*=\operatorname{poly}(d,L,C)$. Tightening $e_*$ by a constant factor
changes only constants and logarithmic factors; it introduces no exponential
dependence on depth. This proves the target-cylinder assertion used above.

\subsection{A trusted-region numerical version}

If the numerical support condition is imposed only on a fixed trusted
region $G$ with $D(G^c)\le\rho$, while the training-margin likelihood is
still imposed on \emph{all} training points, the same proof uses masked
ramp losses $\mathbf1_G(x)\psi(f(x)h(x))$. It gives margin-failure risk
at most $\rho+\varepsilon$ with the same polynomial bound. A finite catalog
of trusted regions can be handled by the same mixture construction as in
Corollary~\ref{cor:trusted-catalog}, at an additional $\log M$ cost.
The global target cylinder belongs to every component.

\subsection{A soft numerical-stability prior with full support}

The hard numerical condition can also be replaced by a soft penalty. Define
\[
 \rho_{\rm num}(w)=\min_x\inf\bigl\{\|\eta\|_*:
       |z_w(x;\eta)-z_w(x)|\ge1/16\bigr\}.
\]
For every finite parameter array, continuity on the finite input cube
gives $\rho_{\rm num}(w)>0$, including when some clean outputs vanish.
Output-only noise gives $\rho_{\rm num}(w)\le1/16$.
For fixed $p,\lambda>0$, use the prior density
\[
 d\nu_{N,t}(w)\ \propto\
       \exp\bigl[-(N/t^2)\lambda\rho_{\rm num}(w)^{-p}\bigr]
       \,d\mu_{N,t}(w),
\]
and retain the training-margin likelihood~\eqref{eq:numeric-likelihood}.
This density is positive at every finite parameter array; it retains the
support of the raw Gaussian and needs no clean-margin exclusion anywhere.

The soft-prior calculation from Theorem~\ref{thm:soft-margin-prior} applies
with $\rho_{\rm num}$ in place of the sign margin. The target cylinder has
$\rho_{\rm num}\ge\tau_0$ and raw mass $e^{-ND_*/t^2}$.
Set $B_*=D_*+\lambda\tau_0^{-p}$. The training partition is at least
$e^{-NB_*/t^2}$; posterior KL is at most $NB_*/t^2$; and
\[
 Q_S\{\rho_{\rm num}^{-p}>u\}
   \le\exp\{(N/t^2)(B_*-\lambda u)\}.
\]
Choosing $u=[B_*+\log(1/\beta)]/\lambda$ discards mass at most $\beta$.
The retained law has scaled KL parameter at most $B_*+1$ and numerical
stability at radius $r=\tfrac12u^{-1/p}\le1/32$.
Lemma~\ref{lem:numeric-compression} and the same ramp-loss argument prove
the conclusion of Theorem~\ref{thm:numeric-prior}, with a polynomial bound
for fixed $p,\lambda$. The prior probabilities have changed; this is not a
theorem for unchanged independent Gaussian weights.
One may also minimize only over a trusted region $G$, obtaining the
corresponding error bound $D(G^c)+\varepsilon$.

This theorem does not identify sign-only conditioning with positive-margin
conditioning. It shows that a global clean-output lower bound is unnecessary
when numerical internal stability is imposed and the observed labels retain
a positive output margin. Numerical stability is not a universal consequence
of ordinary gradient descent: the earlier fragile scalar-chain construction
also violates this condition for sufficiently large depth. For example, at
even $L\ge1024$, its displayed attack bound divided by the present radius
is at most $4000e^{-\sqrt L/4}/\sqrt L<1$, while the attack changes the output
from $+1$ to $-1$. The same strict inequalities persist in its open attraction
neighborhood for gradient descent.
